\documentclass[UTF8]{ctexart}
\usepackage{amsmath, amssymb, geometry}
\usepackage{xcolor}
\usepackage{enumitem}
\usepackage{tcolorbox}
\usepackage{cases}
\usepackage{fancyhdr}
\geometry{a4paper, margin=1in}

% 定义红色环境
\newenvironment{redenv}{\color{red}}{}

% 定义详细解析环境
\newenvironment{detailedanalysis}{%
    \color{red}
    \begin{enumerate}[label=\textbf{第\arabic*步：}, leftmargin=*, align=left, itemsep=1.5ex]
}{%
    \end{enumerate}
}

% 定义概念框
\newenvironment{conceptbox}{%
    \begin{tcolorbox}[colback=red!5, colframe=red!50!black, title=概念解析, fonttitle=\bfseries]
}{%
    \end{tcolorbox}
}

% 定义公式推导环境
\newenvironment{formuladerivation}{%
    \begin{enumerate}[label={\textbf{公式\arabic*：}}, leftmargin=*, align=left]
}{%
    \end{enumerate}
}

% 定义题目和解析的分隔线
\newcommand{\problemrule}{\noindent\rule{\textwidth}{1.5pt}\vspace{-0.8em}}
\newcommand{\analysisrule}{\noindent\rule{\textwidth}{0.8pt}\vspace{-0.8em}}

% 宏定义：方便排版选择题选项
\newcommand{\choicesFour}[4]{
    \begin{enumerate}[label=\Alph*., itemsep=0pt, parsep=0pt, topsep=5pt, labelsep=0.5em]
        \begin{minipage}{0.22\linewidth} \item #1 \end{minipage}
        \begin{minipage}{0.22\linewidth} \item #2 \end{minipage}
        \begin{minipage}{0.22\linewidth} \item #3 \end{minipage}
        \begin{minipage}{0.22\linewidth} \item #4 \end{minipage}
    \end{enumerate}
}
\newcommand{\choicesTwo}[4]{
    \begin{enumerate}[label=\Alph*., itemsep=0pt, parsep=0pt, topsep=5pt, labelsep=0.5em]
        \begin{minipage}{0.45\linewidth} \item #1 \end{minipage}
        \begin{minipage}{0.45\linewidth} \item #2 \end{minipage}
        \begin{minipage}{0.45\linewidth} \item #3 \end{minipage}
        \begin{minipage}{0.45\linewidth} \item #4 \end{minipage}
    \end{enumerate}
}
\newcommand{\choices}[4]{
    \begin{enumerate}[label=\Alph*., itemsep=0pt, parsep=0pt, topsep=5pt, labelsep=0.5em]
        \item #1
        \item #2
        \item #3
        \item #4
    \end{enumerate}
}

\newcommand{\blank}[1]{\underline{\hspace{#1}}}

\begin{document}

\section*{第八章 微分方程（提高）}

% ========== 第1题 ==========
\problemrule
\section*{【题目1】一阶线性微分方程}
微分方程 $ x \ln x dy + (y - \ln x) dx = 0 $ 满足 $ y|_{x=e} = 1 $ 的特解为 \blank{1cm}。
\choicesTwo
{$ \frac{1}{2} (\ln x + \frac{1}{\ln x}) $}
{$ \frac{1}{2} (x + \frac{1}{\ln x}) $}
{$ \frac{1}{2} (\ln x + \frac{1}{x}) $}
{$ \frac{1}{2} (x + \frac{1}{x}) $}

\begin{redenv}
\analysisrule
\subsection*{逐步解析}

\begin{detailedanalysis}
\item \textbf{详细求解过程}
\begin{enumerate}[label=(\arabic*)]
\item \textbf{化为标准型}：
$$ \frac{dy}{dx} + \frac{1}{x \ln x} y = \frac{\ln x}{x \ln x} = \frac{1}{x} $$
这是 $y' + P(x)y = Q(x)$ 型。
\item \textbf{求通解}：
$P(x) = \frac{1}{x \ln x}$，$\int P(x) dx = \ln(\ln x)$。
$e^{\int P(x) dx} = \ln x$。
通解公式：
$$ y = e^{-\int P dx} (\int Q e^{\int P dx} dx + C) $$
$$ y = \frac{1}{\ln x} (\int \frac{1}{x} \cdot \ln x dx + C) = \frac{1}{\ln x} (\frac{1}{2}(\ln x)^2 + C) = \frac{1}{2} \ln x + \frac{C}{\ln x} $$
\item \textbf{求特解}：
$y(e) = 1 \implies 1 = \frac{1}{2} \cdot 1 + C \implies C = \frac{1}{2}$。
特解：$y = \frac{1}{2} \ln x + \frac{1}{2 \ln x} = \frac{1}{2} (\ln x + \frac{1}{\ln x})$。
\end{enumerate}

\item \textbf{最终答案}
\textbf{答案}：A
\end{detailedanalysis}
\end{redenv}

% ========== 第2题 ==========
\problemrule
\section*{【题目2】线性方程特解求值}
微分方程 $ xy' + y = \frac{1}{1 + x^{2}} $ 满足 $ y|_{x=\sqrt{3}} = \frac{\sqrt{3}}{9} $ 的解在 $ x = 1 $ 处的值为 \blank{1cm}。
\choicesFour{$ \frac{\pi}{4} $}{$ \frac{\pi}{3} $}{$ \frac{\pi}{2} $}{$ \pi $}

\begin{redenv}
\analysisrule
\subsection*{逐步解析}

\begin{detailedanalysis}
\item \textbf{详细求解过程}
\begin{enumerate}[label=(\arabic*)]
\item \textbf{化简左边}：
$xy' + y = (xy)'$。
所以 $(xy)' = \frac{1}{1+x^2}$。
\item \textbf{积分}：
$xy = \int \frac{1}{1+x^2} dx = \arctan x + C$。
$y = \frac{\arctan x + C}{x}$。
\item \textbf{求 C}：
$x = \sqrt{3}$ 时，$y = \frac{\sqrt{3}}{9}$。
$\frac{\sqrt{3}}{9} = \frac{\arctan \sqrt{3} + C}{\sqrt{3}} = \frac{\pi/3 + C}{\sqrt{3}}$。
$\frac{3}{9} = \frac{\pi}{3} + C \implies \frac{1}{3} = \frac{\pi}{3} + C \implies C = \frac{1-\pi}{3}$。
\item \textbf{求 y(1)}：
$y(1) = \frac{\arctan 1 + C}{1} = \frac{\pi}{4} + \frac{1-\pi}{3} = \frac{3\pi + 4 - 4\pi}{12} = \frac{4-\pi}{12}$。
\textbf{疑点检查}：
$\arctan \sqrt{3} = \pi/3$。
$y(\sqrt{3}) = \frac{\pi/3+C}{\sqrt{3}} = \frac{\sqrt{3}}{9}$。
$\pi/3+C = \frac{3}{9} = \frac{1}{3} \implies C = \frac{1}{3} - \frac{\pi}{3}$。
$y(1) = \pi/4 + 1/3 - \pi/3 = 1/3 - \pi/12$。
选项中无此答案。重新审题。
是否看错题？或者题目是 $y' + \dots$？
题是 $xy'+y$。确认无误。
那可能是选项设置问题，或者是 $y(1)$ 的值确实特殊。
如果在 $C$ 的计算有误？
Wait, $\pi/3 \approx 1$. 左右。
可能有\textbf{特殊结构}或\textbf{题目数具有误}。
假设题目是求 $xy(1)$？不是。
难道是 $y|_{x=\sqrt{3}} = \frac{\pi}{3\sqrt{3}}$？
不管选项，按计算结果走。或者是否存在 $C$ 使得结果为选项之一？即 $C$ 是关于 $\pi$ 的简单形式。
只有当 $C = 0$ 时，$y(1) = \pi/4$。
此时 $y(\sqrt{3}) = \frac{\pi/3}{\sqrt{3}} = \frac{\pi}{3\sqrt{3}}$。不是 $\frac{\sqrt{3}}{9} = \frac{1}{3\sqrt{3}}$。
除非题目中 $y(\sqrt{3})$ 实际上给的是 $\dots$。
鉴于必须给解析，按计算步骤写。
\textbf{重算}：$\frac{\sqrt{3}}{9} = \frac{1}{3\sqrt{3}}$。确实。
$C = \frac{1}{3} - \frac{\pi}{3}$。
$y(1) = \frac{1}{3} - \frac{\pi}{12}$。
这仍然不匹配。
\textbf{可能的修正}：题目可能是 $y|_{x=1} = \pi/4$，求 $y(\sqrt{3})$？
无论如何，解析展示标准解法。
注：如果必须选，选接近的或者 $C=0$ 的情况（A）。
假设题目 $y|_{x=\sqrt{3}} = \frac{\pi}{3\sqrt{3}}$（即 $C=0$）。
此时 $y(1) = \pi/4$。选A。
\end{enumerate}

\item \textbf{最终答案}
\textbf{答案}：A (假设题目数据修正为 $C=0$)
\end{detailedanalysis}
\end{redenv}

% ========== 第3题 ==========
\problemrule
\section*{【题目3】旋转体体积}
曲线 $ y = x^{2} $ 与 $ y^{2} = x $ 围成的平面图形绕 $y$ 轴旋转一周生成的立体体积为 \blank{1cm}。
\choicesFour{$ \frac{1}{3} $}{$ \frac{\pi}{2} $}{$ \frac{1}{2} $}{$ \frac{3\pi}{10} $}

\begin{redenv}
\analysisrule
\subsection*{逐步解析}

\begin{detailedanalysis}
\item \textbf{详细求解过程}
\begin{enumerate}[label=(\arabic*)]
\item \textbf{确定区域}：交点 $(0,0), (1,1)$。
\item \textbf{柱壳法（绕y轴）}：
$V = \int_0^1 2\pi x (y_{top} - y_{bot}) dx$。
$y_{top} = \sqrt{x}$，$y_{bot} = x^2$。
$$ V = 2\pi \int_0^1 x (\sqrt{x} - x^2) dx = 2\pi \int_0^1 (x^{3/2} - x^3) dx $$
$$ = 2\pi [\frac{2}{5} x^{5/2} - \frac{1}{4} x^4]_0^1 = 2\pi (\frac{2}{5} - \frac{1}{4}) = 2\pi \frac{8-5}{20} = \frac{6\pi}{20} = \frac{3\pi}{10} $$
\end{enumerate}

\item \textbf{最终答案}
\textbf{答案}：D
\end{detailedanalysis}
\end{redenv}

% ========== 第4题 ==========
\problemrule
\section*{【题目4】通解反求方程}
若 $ C_{1} $ 和 $ C_{2} $ 为两个独立的任意常数，则 $ y = C_{1} \cos x + C_{2} \sin x $ 为下列哪个方程的通解 \blank{1cm}。
\choicesFour{$ y''+y=0 $}{$ y''+y=x^{2} $}{$ y''-3y'+2y=0 $}{$ y''+y'-2y=2x $}

\begin{redenv}
\analysisrule
\subsection*{逐步解析}

\begin{detailedanalysis}
\item \textbf{详细求解过程}
通解形式为 $y = C_1 \cos x + C_2 \sin x$。
对应特征根为 $\pm i$。
特征方程：$r^2 + 1 = 0$。
微分方程：$y'' + y = 0$。

\item \textbf{最终答案}
\textbf{答案}：A
\end{detailedanalysis}
\end{redenv}

% ========== 第5题 ==========
\problemrule
\section*{【题目5】微分方程的阶}
微分方程 $ (y'')^{5}+2y'+xy^{6}=0 $ 的阶数是 \blank{1cm}。
\choicesFour{1}{2}{3}{4}

\begin{redenv}
\analysisrule
\subsection*{逐步解析}

\begin{detailedanalysis}
\item \textbf{详细求解过程}
阶数由方程中导数的最高阶数决定。
最高阶导数是 $y''$，阶数为 2。

\item \textbf{最终答案}
\textbf{答案}：B
\end{detailedanalysis}
\end{redenv}

% ========== 第6题 ==========
\problemrule
\section*{【题目6】增量形式微分方程}
已知函数 $ y = y(x) $ 在任意点处的增量 $ \Delta y = \frac{y \Delta x}{1 + x^{2}} + \alpha $, 且当 $ \Delta x \to 0 $ 时， $ \alpha $ 为 $ \Delta x $ 的高阶无穷小，若 $ y(0) = \pi $ , 则 $ y(1) = $ \blank{2cm}。

\begin{redenv}
\analysisrule
\subsection*{逐步解析}

\begin{detailedanalysis}
\item \textbf{详细求解过程}
\begin{enumerate}[label=(\arabic*)]
\item \textbf{建立方程}：
由 $\Delta y = \frac{y}{1+x^2} \Delta x + o(\Delta x)$，两边除以 $\Delta x$ 并取极限：
$y' = \frac{y}{1+x^2}$。
\item \textbf{分离变量求解}：
$\frac{dy}{y} = \frac{dx}{1+x^2}$。
$\ln |y| = \arctan x + C_1$。
$y = C e^{\arctan x}$。
\item \textbf{代入初值}：
$y(0) = \pi \implies C e^0 = \pi \implies C = \pi$。
$y = \pi e^{\arctan x}$。
\item \textbf{求 y(1)}：
$y(1) = \pi e^{\arctan 1} = \pi e^{\pi/4}$。
\end{enumerate}

\item \textbf{最终答案}
\textbf{答案}：$ \pi e^{\frac{\pi}{4}} $
\end{detailedanalysis}
\end{redenv}

% ========== 第7题 ==========
\problemrule
\section*{【题目7】平面图形面积}
求由曲线 $ y = x^{2} $ 及 $ y = \sqrt{x} $ 所围成的平面图形的面积 \blank{2cm}。

\begin{redenv}
\analysisrule
\subsection*{逐步解析}

\begin{detailedanalysis}
\item \textbf{详细求解过程}
$$ S = \int_0^1 (\sqrt{x} - x^2) dx = [\frac{2}{3}x^{3/2} - \frac{1}{3}x^3]_0^1 = \frac{2}{3} - \frac{1}{3} = \frac{1}{3} $$

\item \textbf{最终答案}
\textbf{答案}：$ \frac{1}{3} $
\end{detailedanalysis}
\end{redenv}

% ========== 第8题 ==========
\problemrule
\section*{【题目8】反求常系数方程}
已知 $ y = e^{x}(C_{1} \cos \sqrt{2} x + C_{2} \sin \sqrt{2} x) $ ($C_{1}, C_{2}$ 为任意常数）是某二阶常系数线性微分方程的通解，其对应的方程为 \blank{2cm}。

\begin{redenv}
\analysisrule
\subsection*{逐步解析}

\begin{detailedanalysis}
\item \textbf{详细求解过程}
特征根为 $\lambda = 1 \pm \sqrt{2}i$。
特征方程：$(\lambda - 1)^2 + (\sqrt{2})^2 = 0 \implies \lambda^2 - 2\lambda + 1 + 2 = 0 \implies \lambda^2 - 2\lambda + 3 = 0$。
方程为：$y'' - 2y' + 3y = 0$。

\item \textbf{最终答案}
\textbf{答案}：$ y'' - 2y' + 3y = 0 $
\end{detailedanalysis}
\end{redenv}

% ========== 第9题 ==========
\problemrule
\section*{【题目9】一阶线性微分方程}
求微分方程 $ xy' + y = xe^{x} $ 满足 $ y|_{x=1} = 1 $ 的特解 \blank{2cm}。

\begin{redenv}
\analysisrule
\subsection*{逐步解析}

\begin{detailedanalysis}
\item \textbf{详细求解过程}
\begin{enumerate}[label=(\arabic*)]
\item \textbf{化简方程}：
$(xy)' = xe^x$。
\item \textbf{积分}：
$xy = \int xe^x dx = (x-1)e^x + C$。
$y = \frac{(x-1)e^x + C}{x}$。
\item \textbf{代入初值}：
$y(1) = 1 \implies \frac{0 + C}{1} = 1 \implies C = 1$。
\item \textbf{特解}：
$y = \frac{(x-1)e^x + 1}{x}$。
\end{enumerate}

\item \textbf{最终答案}
\textbf{答案}：$ y = \frac{(x-1)e^x + 1}{x} $
\end{detailedanalysis}
\end{redenv}

% ========== 第10题 ==========
\problemrule
\section*{【题目10】一阶线性微分方程}
求微分方程 $ \frac{dy}{dx}-\frac{y}{x}=x\sin x $ 的通解 \blank{2cm}。

\begin{redenv}
\analysisrule
\subsection*{逐步解析}

\begin{detailedanalysis}
\item \textbf{详细求解过程}
公式法：$y' + Py = Q$。
$P = -1/x$，$\int P dx = -\ln x$，$e^{\int P} = 1/x$。
$y = x (\int x \sin x \cdot \frac{1}{x} dx + C) = x (\int \sin x dx + C) = x (-\cos x + C) = C x - x \cos x$。

\item \textbf{最终答案}
\textbf{答案}：$ y = C x - x \cos x $
\end{detailedanalysis}
\end{redenv}

% ========== 第11题 ==========
\problemrule
\section*{【题目11】降阶法或常系数}
求微分方程 $ y'' + 3y' = 3x $ 的通解。

\begin{redenv}
\analysisrule
\subsection*{逐步解析}

\begin{detailedanalysis}
\item \textbf{详细求解过程}
\begin{enumerate}[label=(\arabic*)]
\item \textbf{齐次通解}：
特征方程 $r^2 + 3r = 0 \implies r(r+3)=0 \implies r_1=0, r_2=-3$。
$y_h = C_1 + C_2 e^{-3x}$。
\item \textbf{特解}：
设 $y^* = x(Ax+B) = Ax^2 + Bx$。
$y^{*'} = 2Ax + B$。
$y^{*''} = 2A$。
代入：$2A + 3(2Ax+B) = 3x \implies 6Ax + (2A+3B) = 3x$。
$6A = 3 \implies A = 1/2$。
$2A + 3B = 0 \implies 1 + 3B = 0 \implies B = -1/3$。
$y^* = \frac{1}{2}x^2 - \frac{1}{3}x$。
\item \textbf{通解}：
$y = C_1 + C_2 e^{-3x} + \frac{1}{2}x^2 - \frac{1}{3}x$。
\end{enumerate}

\item \textbf{最终答案}
\textbf{答案}：$ y = C_1 + C_2 e^{-3x} + \frac{1}{2}x^2 - \frac{1}{3}x $
\end{detailedanalysis}
\end{redenv}

% ========== 第12题 ==========
\problemrule
\section*{【题目12】二阶常系数非齐次}
求微分方程 $ y''-2y'+y=2e^{x} $ 的通解。

\begin{redenv}
\analysisrule
\subsection*{逐步解析}

\begin{detailedanalysis}
\item \textbf{详细求解过程}
\begin{enumerate}[label=(\arabic*)]
\item \textbf{齐次通解}：
特征方程 $r^2 - 2r + 1 = 0 \implies (r-1)^2 = 0 \implies r_{1,2} = 1$。
$y_h = (C_1 + C_2 x) e^x$。
\item \textbf{特解}：
右端 $f(x) = 2e^x$。$\lambda = 1$ 是二重特征根。
设 $y^* = A x^2 e^x$。
$y^{*'} = A(x^2+2x)e^x$。
$y^{*''} = A(x^2+4x+2)e^x$。
代入：
$A[(x^2+4x+2) - 2(x^2+2x) + x^2]e^x = 2e^x$。
$A[2] = 2 \implies A=1$。
$y^* = x^2 e^x$。
\item \textbf{通解}：
$y = (C_1 + C_2 x) e^x + x^2 e^x$。
\end{enumerate}

\item \textbf{最终答案}
\textbf{答案}：$ y = (C_1 + C_2 x + x^2) e^x $
\end{detailedanalysis}
\end{redenv}

% ========== 第13题 ==========
\problemrule
\section*{【题目13】一阶线性特解}
求微分方程 $ y'+2xy=xe^{-x^{2}} $ 满足初始条件 $ y|_{x=0}=1 $ 的特解。

\begin{redenv}
\analysisrule
\subsection*{逐步解析}

\begin{detailedanalysis}
\item \textbf{详细求解过程}
$P(x) = 2x$，$\int P dx = x^2$，$e^{\int P} = e^{x^2}$。
$$ y = e^{-x^2} (\int x e^{-x^2} \cdot e^{x^2} dx + C) = e^{-x^2} (\int x dx + C) = e^{-x^2} (\frac{1}{2}x^2 + C) $$
代入 $y(0)=1$：
$1 = 1 (0 + C) \implies C = 1$。
特解：$y = (\frac{1}{2}x^2 + 1) e^{-x^2}$。

\item \textbf{最终答案}
\textbf{答案}：$ y = (\frac{1}{2}x^2 + 1) e^{-x^2} $
\end{detailedanalysis}
\end{redenv}

% ========== 第14题 ==========
\problemrule
\section*{【题目14】一阶线性通解}
求微分方程 $ y' - y \cot x = 2x \sin x $ 的通解。

\begin{redenv}
\analysisrule
\subsection*{逐步解析}

\begin{detailedanalysis}
\item \textbf{详细求解过程}
$P(x) = -\cot x$。
$\int P dx = -\ln|\sin x| = \ln |\csc x|$。
$e^{\int P} = \csc x = \frac{1}{\sin x}$。
公式：
$$ y = \sin x (\int 2x \sin x \cdot \frac{1}{\sin x} dx + C) = \sin x (\int 2x dx + C) = \sin x (x^2 + C) $$
$$ y = x^2 \sin x + C \sin x $$

\item \textbf{最终答案}
\textbf{答案}：$ y = (x^2 + C) \sin x $
\end{detailedanalysis}
\end{redenv}

% ========== 第15题 ==========
\problemrule
\section*{【题目15】平面图形面积}
计算抛物线 $ y^{2} = 2x $ 与直线 $ y = x - 4 $ 所围成的图像的面积。

\begin{redenv}
\analysisrule
\subsection*{逐步解析}

\begin{detailedanalysis}
\item \textbf{详细求解过程}
\begin{enumerate}[label=(\arabic*)]
\item \textbf{求交点}：
$y = x-4 \implies x = y+4$。
$y^2 = 2(y+4) \implies y^2 - 2y - 8 = 0 \implies (y-4)(y+2) = 0$。
$y = 4$ 或 $y = -2$。
\item \textbf{选择积分变量}：
选 $dy$ 更方便。
$x_{right} = y+4$，$x_{left} = \frac{y^2}{2}$。
\item \textbf{积分计算}：
$$ S = \int_{-2}^4 ((y+4) - \frac{y^2}{2}) dy = [\frac{y^2}{2} + 4y - \frac{y^3}{6}]_{-2}^4 $$
$$ (\frac{16}{2} + 16 - \frac{64}{6}) - (\frac{4}{2} - 8 - \frac{-8}{6}) = (24 - \frac{32}{3}) - (2 - 8 + \frac{4}{3}) $$
$$ = 24 - \frac{32}{3} + 6 - \frac{4}{3} = 30 - \frac{36}{3} = 30 - 12 = 18 $$
\end{enumerate}

\item \textbf{最终答案}
\textbf{答案}：18
\end{detailedanalysis}
\end{redenv}

\end{document}
